\chapter{CONCLUSION}\label{chap_conclusao}

This thesis investigated the current performance of a quantum computer on the task of classifying data from the Iris dataset by comparing it with a classical classifier. The main objective was to evaluate the capability of a real quantum computer on a machine learning problem, specifically in the area of data classification. The Iris dataset was selected as the case study because it is widely used in the scientific community as a benchmark for classification algorithms.

To achieve this objective, we used a supervised quantum machine learning method, the variational quantum classifier, executed on a quantum computer. The quantum classifier was constructed from quantum circuits for encoding classical information into quantum states and for learning. Training consisted of learning the quantum-gate parameters that optimize the chosen cost function.

Through detailed comparisons between the classical classifier and the real quantum computer, we analyzed several aspects, including performance metrics and training-set sizes. In general, the classical classifier outperformed the quantum classifier for most metrics across all evaluated training sizes.

These results suggest that the classical model may be more suitable for the Iris dataset used in this study. Nevertheless, the quantum classifier showed consistent performance across all metrics and training sizes.

Although the quantum classifier produced results below those of the classical classifier, it is important to consider that quantum computing is still an emerging field. Given the consistency and robustness demonstrated by the quantum classifier in this study, further optimization and development in this area may make it possible to match or even surpass the performance of the classical classifier.

This work contributed to the understanding of the current state of quantum computing in relation to machine learning. By analyzing the capability of a real quantum computer on a data classification task, the study identified its limitations and obtained valuable insights for future research and development in this constantly evolving area.

In addition, this study opens the way for future research, particularly research aimed at optimizing quantum classifiers and exploring their potential advantages. Extending this analysis to other datasets and scenarios may provide further insights into the applicability and effectiveness of quantum classifiers.